\subsection{Primeira passagem: o instante em que o FS cruza 1}\label{sec:cruza}
O instante $t_1$ em que o FS cruza 1 pela primeira vez (interpolado linearmente entre as duas horas calculadas
vizinhas; triângulos nas Figs.~\ref{fig:f10} a \ref{fig:f100}) é a grandeza que importa para o Monte Carlo, mais que o
FS mínimo:
\begin{itemize}
  \item \textbf{É a ruptura.} Em $t_1$ o talude perde o equilíbrio; o que vem depois no cálculo (o FS descendo até o
    mínimo e voltando acima de 1 com a drenagem) usa a geometria intacta e as pressões de um talude que já rompeu, e
    não tem sentido físico. A falha de uma realização é o evento de primeira passagem, FS$(t) < 1$ para algum $t$, e a
    P$_f$ é a probabilidade dele. Ela é a mesma do critério ``FS mínimo $< 1$'' quando os instantes são os mesmos, mas
    o $t_1$ dá também a \emph{distribuição do instante de ruptura}, que o FS mínimo não dá.
  \item \textbf{A realização pode parar em $t_1$.} No f10 com 180~mm/h a ruptura ocorre a 10,7~h: parar ali corta 78\,\%
    do horizonte de 48~h (o transiente e as análises do fator que viriam depois). Nas realizações que não rompem, o
    cálculo vai até o fim do horizonte de qualquer jeito.
  \item \textbf{Precisão de $t_1$.} Com o FS de hora em hora e a interpolação linear, o erro em $t_1$ é de uma fração de
    hora. No Monte Carlo, avaliar o FS de hora em hora custa 49 SRM por realização; a alternativa é avaliar em poucos
    instantes e, no intervalo em que o FS cruza 1, refinar o instante por bissecção (o driver já faz isso para a curva
    crítica de chuva, \texttt{crpc=}).
\end{itemize}

\begin{table}[H]\centering\small
\caption{Primeira passagem: $t_1$ (o FS cruza 1 descendo), o instante em que volta a ficar acima de 1 e o tempo
abaixo de 1 (h; interpolação linear entre as horas calculadas).}\label{tab:cruza}
\begin{tabular}{llccccccccc}
\toprule
 & & \multicolumn{3}{c}{(iii)} & \multicolumn{3}{c}{(ii)} & \multicolumn{3}{c}{(i)}\\
\cmidrule(lr){3-5}\cmidrule(lr){6-8}\cmidrule(lr){9-11}
solo & $I_r$ (mm/h) & $t_1$ (h) & volta (h) & abaixo (h) & $t_1$ (h) & volta (h) & abaixo (h) & $t_1$ (h) & volta (h) & abaixo (h)\\
\midrule
f10 & 36 & não cruza &  &  & não cruza &  &  & não cruza &  & \\
f10 & 90 & não cruza &  &  & não cruza &  &  & não cruza &  & \\
f10 & 180 & 10{,}7 & 26{,}4 & 15{,}7 & 10{,}9 & 26{,}5 & 15{,}5 & 5{,}8 & 25{,}0 & 19{,}2\\
\midrule
f50 & 9 & não cruza &  &  & não cruza &  &  & não cruza &  & \\
f50 & 36 & não cruza &  &  & não cruza &  &  & 15{,}0 & 28{,}1 & 13{,}0\\
f50 & 90 & 19{,}3 & 32{,}3 & 13{,}0 & 20{,}5 & 32{,}0 & 11{,}5 & 11{,}3 & 28{,}3 & 17{,}0\\
\midrule
f100 & 3{,}6 & não cruza &  &  & não cruza &  &  & não cruza &  & \\
f100 & 9 & não cruza &  &  & não cruza &  &  & não cruza &  & \\
f100 & 36 & não cruza &  &  & não cruza &  &  & não cruza &  & \\
\bottomrule
\end{tabular}

\end{table}

\paragraph{Observações (variante iii).} Só dois casos rompem: o f10 com 180~mm/h ($t_1 = 10{,}7$~h, 45\,\% da chuva) e
o f50 com 90~mm/h ($t_1 = 19{,}3$~h, 80\,\% da chuva). Nos dois a ruptura ocorre \emph{durante} a chuva, bem antes do
mínimo (24~h), e o FS fica abaixo de 1 por 13--16~h. O f50 com 36~mm/h passa a 0,034 de romper (FS 1,034 em 24~h) e
não cruza; o f100 não chega perto de 1 com chuvas de 24~h.
